\section{The High Bleeding Risk Population}

A specific subset of patients undergoing PCI is classified as High Bleeding
Risk (HBR), defined by the presence of one or more clinical factors that
substantially elevate the probability of a significant bleeding event.
HBR patients represent a particularly challenging population because:
\begin{itemize}
  \item They derive the greatest harm from prolonged DAPT due to their
    elevated baseline bleeding risk.
  \item They are frequently excluded or underrepresented in standard PCI
    trials, which limits the generalizability of existing evidence.
  \item Their clinical profiles are highly heterogeneous, spanning a wide
    range of comorbidities, frailty levels, and concomitant medications.
\end{itemize}

A patient was eligible for MASTER DAPT if at least one of the following ten
criteria, adapted from the Academic Research Consortium for High Bleeding
Risk (ARC-HBR) consensus, applied (Valgimigli et al., \emph{NEJM} 2021,
Supplementary Appendix, Section 1):
\begin{enumerate}
  \item Indication for oral anticoagulation (OAC) for at least 12 months.
  \item A nonaccess-site bleeding episode requiring medical attention within
    the previous 12 months.
  \item A previous bleeding episode requiring hospitalization, if the
    underlying cause had not been definitively treated.
  \item Age $\geq 75$ years.
  \item A systemic condition associated with increased bleeding risk (e.g.,
    thrombocytopenia $<100{,}000/\mathrm{mm}^3$ or a known coagulation
    disorder).
  \item Documented anemia (hemoglobin $<11\,\mathrm{g/dL}$, or transfusion
    within the previous 4 weeks).
  \item Chronic need for steroids or nonsteroidal anti-inflammatory drugs.
  \item A malignancy (other than skin) considered at high bleeding risk.
  \item Stroke or transient ischemic attack within the previous 6 months.
  \item A PRECISE-DAPT score $\geq 25$.
\end{enumerate}
In the MASTER DAPT trial, the mean number of criteria met per patient was
$2.1 \pm 1.1$, and the most prevalent single criterion was age $\geq 75$
years (met in isolation by 17\% of patients), reflecting a substantially
vulnerable population. Note that chronic kidney disease, although a common
comorbidity in this cohort (reported separately in the baseline
characteristics), is not itself one of the ten enrollment criteria: renal
impairment enters indirectly, through its contribution to anemia, the
PRECISE-DAPT score, or malignancy staging.

\section{The MASTER DAPT Trial Design}

MASTER DAPT was a multicenter, open-label, randomized, noninferiority trial
with hierarchical superiority testing. Key design features include:

\begin{itemize}
  \item \textbf{Enrollment period:} February 28, 2017 to December 5, 2019.
  \item \textbf{Sites:} 140 centers across 30 countries.
  \item \textbf{Screened patients:} 5,204; randomized: 4,579 (88.0\%).
  \item \textbf{Randomization:} 1:1, computer-generated, stratified by trial
    site, history of myocardial infarction in the past 12 months, and
    indication for oral anticoagulation.
  \item \textbf{Follow-up:} 335 days with visits at 60, 150, and 335 days
    post-randomization.
\end{itemize}

\paragraph{Patient characteristics at baseline.}
The population enrolled in MASTER DAPT was elderly and clinically complex:
mean age $76.0$ years, 69.3\% male. Prominent comorbidities included
diabetes (33.6\%), chronic kidney disease (19.1\%), heart failure (18.9\%),
history of cerebrovascular events (12.4\%), and peripheral vascular disease
(10.6\%). Concomitant oral anticoagulation was present in 36.4\% of patients.
These characteristics reflect a real-world HBR population that is
substantially more complex than the patients typically studied in earlier
DAPT trials.

\paragraph{Interventions.}
\begin{itemize}
  \item \textbf{Abbreviated therapy group:} Immediate discontinuation of
    DAPT at randomization (approximately 1 month after the index PCI),
    followed by single antiplatelet therapy. For patients on oral
    anticoagulation, single antiplatelet therapy was continued up to 6 months.
  \item \textbf{Standard therapy group:} Continuation of DAPT for at least
    5 additional months (6 months total from index PCI), or at least 2
    additional months for patients on oral anticoagulation.
\end{itemize}

\paragraph{Primary outcomes.}
Three hierarchically ranked primary outcomes, assessed as cumulative
incidences at 335 days:
\begin{enumerate}
  \item \textbf{NACE} (Net Adverse Clinical Events): composite of death
    from any cause, myocardial infarction, stroke, or major bleeding
    (BARC type 3 or 5).
  \item \textbf{MACCE} (Major Adverse Cardiac or Cerebral Events): composite
    of death from any cause, myocardial infarction, or stroke.
  \item \textbf{Major or CRNM bleeding}: composite of bleeding events of
    BARC type 2, 3, or 5.
\end{enumerate}

\paragraph{Primary results.}
\begin{itemize}
  \item \textbf{NACE:} 7.5\% (abbreviated) vs. 7.7\% (standard); difference
    $-0.23$ pp (95\% CI $-1.80$ to $1.33$; $p < 0.001$ for noninferiority).
  \item \textbf{MACCE:} 6.1\% vs. 5.9\%; difference $0.11$ pp (95\% CI
    $-1.29$ to $1.51$; $p = 0.001$ for noninferiority).
  \item \textbf{Bleeding:} 6.5\% vs. 9.4\%; difference $-2.82$ pp (95\% CI
    $-4.40$ to $-1.24$; $p < 0.001$ for superiority).
\end{itemize}

\begin{figure}[h]
\centering
\includegraphics[width=0.85\textwidth]{masterdapt_flow.png}
\caption{MASTER DAPT patient flow diagram (screening, randomization, follow-up). Reproduced from Valgimigli et al., \emph{NEJM} 2021.}
\label{fig:masterdapt_flow}
\end{figure}

\section{The Clinical Decision Problem}

From a statistical decision-making perspective, the MASTER DAPT question is:
\textit{given the characteristics of a specific patient, which DAPT strategy
— abbreviated or standard — minimizes their expected net harm?} This question
is precisely what the CATE formalizes. The CATE $\tau(x)$ for this setting
can be interpreted as the expected difference in NACE (or MACCE, or bleeding)
between the abbreviated and standard therapy, conditional on a patient having
covariate profile $x$.

A positive $\tau(x)$ would indicate that abbreviated therapy leads to a
\textit{higher} event rate than standard therapy for patients with profile $x$
(i.e., standard therapy is preferred), while a negative value indicates
the opposite. A CATE near zero suggests no differential benefit of either
strategy for that patient type.